\chapter{结论}{Conclusion}

\section{工作总结}{Summary}

本文档展示了厦门大学本科毕业论文模板的基本使用方式，覆盖封面信息、中英文摘要、双语目录、章节组织、图表、公式、算法、实验表格、参考文献和示例 PDF。正式论文的结论部分应回扣研究问题，概括方法贡献与实验发现，并避免引入正文中未讨论的新结果。

一个稳妥的总结结构是：首先重述本文关注的问题；其次概括本文采用的方法或实现路径；然后总结主要实验观察；最后说明这些观察对研究问题意味着什么。总结不应简单复制摘要，而应从完整论文的证据出发，给出更凝练的回答。

\section{主要发现}{Findings}

主要发现应来自实验和分析，而不是来自主观判断。例如，可以写“在统一评价协议下，方法 A 在指标 B 上优于基线 C”，也可以写“消融结果显示组件 D 对最终性能有稳定贡献”。如果结果存在波动，应说明波动范围和可能原因。

对于本科毕业论文，清楚呈现实验事实比追求夸张贡献更重要。若方法没有取得预期提升，也可以通过负结果分析、复现经验和误差分析体现工作价值。

\section{局限与展望}{Limitations and Future Work}

局限性应具体、诚实且可验证，例如数据规模有限、评价协议不完整、实验资源受限、方法假设较强或缺少真实场景验证。未来工作可以围绕这些局限提出自然延伸，但不宜写成泛泛的方向罗列。

展望部分可以从三个层面展开：数据层面扩展更大或更真实的数据；方法层面改进模型结构或训练目标；应用层面验证方法在实际系统中的可用性。每个方向都应与本文工作自然衔接。
